\chapter{LX-019 --- Localization Port Contract}\label{record-195}

\noindent{\footnotesize\color{BookMuted}\path{docs/design/LX-019-localization-port-contract.md}\par}\smallskip

Legend: \href{https://github.com/flyingrobots/bijou/blob/1853ad902f6bf219326409e7f743a5a1d6517dbc/docs/legends/LX-localization-and-bidirectionality.md}{LX --- Localization and Bidirectionality}

\section{Why This Cycle Exists}\label{record-195-why-this-cycle-exists}

Bijou already has a runtime catalog engine, formatter ports, catalog loaders, string-table tooling, and DOGFOOD locale preference ports. Those pieces make localization possible, but DOGFOOD still tends to pass the concrete i18n runtime through rendering code and ask it for naked strings.

That is too low-level for the long-term architecture.

Views, blocks, and docs surfaces should ask an application-facing localization port for a localized object:

\begin{CodeBlock}
const title = localization.resolve({
  key: { namespace: 'bijou.dogfood', id: 'blocks.preview.title' },
});
\end{CodeBlock}

The result should preserve runtime truth:

\begin{itemize}
\tightlist
\item
  the requested key
\item
  selected locale
\item
  direction
\item
  entry kind
\item
  translated / fallback / missing status
\item
  localized value
\item
  issues and facts for tooling
\end{itemize}

The view should not know how catalogs were loaded, where JSON files live, which CSV source produced them, or how fallback catalogs were attached.

\section{Core Rule}\label{record-195-core-rule}

Localization in Bijou is a port/adapter boundary.

\begin{itemize}
\tightlist
\item
  Ports live where the localization vocabulary lives.
\item
  Adapters live where the external dependency lives.
\item
  Composition lives at the app boundary.
\item
  Views receive localized objects, not filesystem, catalog, or runtime machinery.
\end{itemize}

\section{Boundary}\label{record-195-boundary}

This cycle is about the application-facing localization port contract and DOGFOOD wiring.

It is not about:

\begin{itemize}
\tightlist
\item
  expanding translation coverage
\item
  adding a localization dashboard
\item
  adding a translation workbench
\item
  changing the CSV string-table workflow
\item
  loading every locale at runtime
\item
  implementing remote localization services
\item
  making blocks fetch localization data directly
\item
  replacing the lower-level i18n runtime
\end{itemize}

\section{Vocabulary}\label{record-195-vocabulary}

\subsection{Localization Port}\label{record-195-localization-port}

The app-facing port that resolves a localization request into a structured localized object.

It should answer:

\begin{CodeBlock}
Given this key and interpolation values, what localized object should this view
render in the selected locale?
\end{CodeBlock}

It should not answer:

\begin{CodeBlock}
Where are the catalog files?
How do I read the filesystem?
How do I compile CSV rows?
How do I persist locale preferences?
\end{CodeBlock}

\subsection{Localized Object}\label{record-195-localized-object}

The immutable result returned by the localization port.

It contains:

\begin{itemize}
\tightlist
\item
  key
\item
  locale
\item
  direction
\item
  fallback locale
\item
  entry kind
\item
  resolution status
\item
  value
\item
  issues
\item
  facts
\end{itemize}

The value may be a message string, data payload, or resource payload depending on the catalog entry kind.

\subsection{Runtime Adapter}\label{record-195-runtime-adapter}

An adapter that wraps \texttt{I\allowbreak{}1\allowbreak{}8\allowbreak{}n\allowbreak{}R\allowbreak{}u\allowbreak{}n\allowbreak{}t\allowbreak{}i\allowbreak{}m\allowbreak{}e} and exposes it through \texttt{L\allowbreak{}o\allowbreak{}c\allowbreak{}a\allowbreak{}l\allowbreak{}i\allowbreak{}z\allowbreak{}a\allowbreak{}t\allowbreak{}i\allowbreak{}o\allowbreak{}n\allowbreak{}P\allowbreak{}o\allowbreak{}r\allowbreak{}t}.

The runtime continues to own catalog state, fallback behavior, locale switching, references, and formatting. The port makes that behavior consumable by application code without leaking the runtime surface everywhere.

Runtime convenience methods that remain public, such as \texttt{t\allowbreak{}(\allowbreak{})}, should behave like stable API functions instead of receiver-sensitive object methods. Callers may pass or destructure these helpers without preserving a JavaScript \texttt{t\allowbreak{}h\allowbreak{}i\allowbreak{}s} binding; the runtime state remains closure-owned at the adapter boundary.

\subsection{Node / Filesystem Adapters}\label{record-195-node-filesystem-adapters}

Adapters that read selected-locale generated JSON catalogs from disk belong in \texttt{@\allowbreak{}f\allowbreak{}l\allowbreak{}y\allowbreak{}i\allowbreak{}n\allowbreak{}g\allowbreak{}r\allowbreak{}o\allowbreak{}b\allowbreak{}o\allowbreak{}t\allowbreak{}s\allowbreak{}/\allowbreak{}b\allowbreak{}i\allowbreak{}j\allowbreak{}o\allowbreak{}u\allowbreak{}-\allowbreak{}i\allowbreak{}1\allowbreak{}8\allowbreak{}n\allowbreak{}-\allowbreak{}t\allowbreak{}o\allowbreak{}o\allowbreak{}l\allowbreak{}s\allowbreak{}-\allowbreak{}n\allowbreak{}o\allowbreak{}d\allowbreak{}e} or application composition code.

They should return catalogs or construct runtime-backed localization ports. They should not redefine the generic localization port contract.

\section{Package Split}\label{record-195-package-split}

The contract belongs in \texttt{@\allowbreak{}f\allowbreak{}l\allowbreak{}y\allowbreak{}i\allowbreak{}n\allowbreak{}g\allowbreak{}r\allowbreak{}o\allowbreak{}b\allowbreak{}o\allowbreak{}t\allowbreak{}s\allowbreak{}/\allowbreak{}b\allowbreak{}i\allowbreak{}j\allowbreak{}o\allowbreak{}u\allowbreak{}-\allowbreak{}i\allowbreak{}1\allowbreak{}8\allowbreak{}n}:

\begin{CodeBlock}
packages/bijou-i18n
  LocalizationPort
  LocalizationRequest
  LocalizedObject
  LocalizationStatus
  LocalizationIssue
  LocalizationFact
  createRuntimeLocalizationPort()
\end{CodeBlock}

Node/file adapters stay in \texttt{@\allowbreak{}f\allowbreak{}l\allowbreak{}y\allowbreak{}i\allowbreak{}n\allowbreak{}g\allowbreak{}r\allowbreak{}o\allowbreak{}b\allowbreak{}o\allowbreak{}t\allowbreak{}s\allowbreak{}/\allowbreak{}b\allowbreak{}i\allowbreak{}j\allowbreak{}o\allowbreak{}u\allowbreak{}-\allowbreak{}i\allowbreak{}1\allowbreak{}8\allowbreak{}n\allowbreak{}-\allowbreak{}t\allowbreak{}o\allowbreak{}o\allowbreak{}l\allowbreak{}s\allowbreak{}-\allowbreak{}n\allowbreak{}o\allowbreak{}d\allowbreak{}e}:

\begin{CodeBlock}
packages/bijou-i18n-tools-node
  CSV and TSV string-table files
  generated runtime catalog JSON files
  selected-locale catalog directory loaders
\end{CodeBlock}

DOGFOOD composition lives in \texttt{e\allowbreak{}x\allowbreak{}a\allowbreak{}m\allowbreak{}p\allowbreak{}l\allowbreak{}e\allowbreak{}s\allowbreak{}/\allowbreak{}d\allowbreak{}o\allowbreak{}c\allowbreak{}s}:

\begin{CodeBlock}
examples/docs
  selected locale resolution
  English fallback catalog attachment
  development missing-marker policy
  runtime-backed LocalizationPort assembly
\end{CodeBlock}

Views and block previews should receive the port or localized objects. They should not receive filesystem loaders, CSV data, generated catalog paths, or provider handles.

\section{Expected Flow}\label{record-195-expected-flow}

% Presentation of original arrow notation. See flow-manifest.json.
\begin{center}
\begin{tikzpicture}[book flow node/.style={draw=BookHeading,line width=.65pt,fill=BookPanel,rounded corners=1.5pt,align=center,inner sep=5pt,minimum height=10mm,font=\fontsize{9}{11}\selectfont,execute at begin node={\hyphenpenalty=10000\exhyphenpenalty=10000}},book flow edge/.style={-{Latex[length=2mm]},draw=BookHeading,line width=.7pt}]
\node[book flow node,text width=96mm] (n0) at (0,0.0) {CSV source string table};
\node[book flow node,text width=96mm] (n1) at (0,-1.55) {generated per-locale JSON catalogs};
\node[book flow node,text width=96mm] (n2) at (0,-3.1) {selected-locale catalog loader adapter};
\node[book flow node,text width=96mm] (n3) at (0,-4.65) {I18n\allowbreak{}Runtime};
\node[book flow node,text width=96mm] (n4) at (0,-6.2) {Localization\allowbreak{}Port};
\node[book flow node,text width=96mm] (n5) at (0,-7.75) {Localized\allowbreak{}Object};
\node[book flow node,text width=96mm] (n6) at (0,-9.3) {DOGFOOD /\allowbreak{} Blocks /\allowbreak{} shell surfaces};
\draw[book flow edge] (n0) -- (n1);
\draw[book flow edge] (n1) -- (n2);
\draw[book flow edge] (n2) -- (n3);
\draw[book flow edge] (n3) -- (n4);
\draw[book flow edge] (n4) -- (n5);
\draw[book flow edge] (n5) -- (n6);
\end{tikzpicture}
\end{center}


The fallback posture remains:

% Presentation of original arrow notation. See flow-manifest.json.
\begin{center}
\begin{tikzpicture}[book flow node/.style={draw=BookHeading,line width=.65pt,fill=BookPanel,rounded corners=1.5pt,align=center,inner sep=5pt,minimum height=10mm,font=\fontsize{9}{11}\selectfont,execute at begin node={\hyphenpenalty=10000\exhyphenpenalty=10000}},book flow edge/.style={-{Latex[length=2mm]},draw=BookHeading,line width=.7pt}]
\node[book flow node,text width=48mm] (n0) at (-3,0) {production};
\node[book flow node,text width=48mm] (n1) at (-3,-1.8) {selected locale value};
\node[book flow node,text width=48mm] (n2) at (-3,-3.6) {English fallback catalog};
\node[book flow node,text width=48mm] (n3) at (3,0) {non-production};
\node[book flow node,text width=48mm] (n4) at (3,-1.8) {selected locale value};
\node[book flow node,text width=48mm] (n5) at (3,-3.6) {loud missing-localization object};
\node[book flow node,text width=48mm] (n6) at (3,-5.4) {English fallback only when missing markers are disabled};
\draw[book flow edge] (n0) -- (n1);
\draw[book flow edge] (n1) -- (n2);
\draw[book flow edge] (n3) -- (n4);
\draw[book flow edge] (n4) -- (n5);
\draw[book flow edge] (n5) -- (n6);
\end{tikzpicture}
\end{center}


\section{DOGFOOD Implications}\label{record-195-dogfood-implications}

DOGFOOD should keep the concrete runtime at the composition boundary because the frame shell still consumes \texttt{I\allowbreak{}1\allowbreak{}8\allowbreak{}n\allowbreak{}R\allowbreak{}u\allowbreak{}n\allowbreak{}t\allowbreak{}i\allowbreak{}m\allowbreak{}e} and the app still needs locale activation. Ordinary DOGFOOD text lookup should move behind \texttt{L\allowbreak{}o\allowbreak{}c\allowbreak{}a\allowbreak{}l\allowbreak{}i\allowbreak{}z\allowbreak{}a\allowbreak{}t\allowbreak{}i\allowbreak{}o\allowbreak{}n\allowbreak{}P\allowbreak{}o\allowbreak{}r\allowbreak{}t\allowbreak{}.\allowbreak{}r\allowbreak{}e\allowbreak{}s\allowbreak{}o\allowbreak{}l\allowbreak{}v\allowbreak{}e\allowbreak{}(\allowbreak{})}.

The app may still hold both:

\begin{CodeBlock}
const runtime = createI18nRuntime(...);
const localization = createRuntimeLocalizationPort(runtime);
\end{CodeBlock}

But render helpers should prefer:

\begin{CodeBlock}
localization.resolve({
  key: { namespace: 'bijou.dogfood', id: 'settings.language.label' },
});
\end{CodeBlock}

over:

\begin{CodeBlock}
runtime.t({ namespace: 'bijou.dogfood', id: 'settings.language.label' });
\end{CodeBlock}

\section{Acceptance Criteria}\label{record-195-acceptance-criteria}

\begin{itemize}
\tightlist
\item
  \texttt{@\allowbreak{}f\allowbreak{}l\allowbreak{}y\allowbreak{}i\allowbreak{}n\allowbreak{}g\allowbreak{}r\allowbreak{}o\allowbreak{}b\allowbreak{}o\allowbreak{}t\allowbreak{}s\allowbreak{}/\allowbreak{}b\allowbreak{}i\allowbreak{}j\allowbreak{}o\allowbreak{}u\allowbreak{}-\allowbreak{}i\allowbreak{}1\allowbreak{}8\allowbreak{}n} exports a first-class localization port contract.
\item
  Runtime lookup helpers stay receiver-independent so app code can pass them around without binding \texttt{t\allowbreak{}h\allowbreak{}i\allowbreak{}s}.
\item
  The runtime adapter returns immutable localized objects.
\item
  Localized values enforce the portable JSON-shaped catalog boundary with deterministic errors for cyclic graphs, symbol-keyed properties, non-enumerable properties, accessors, class instances, functions, symbols, and bigint values.
\item
  Localized objects preserve key, locale, direction, kind, status, value, issues, and facts.
\item
  Missing selected-locale translations remain inspectable as missing localization results in non-production posture.
\item
  Production fallback remains explicit through fallback catalogs.
\item
  DOGFOOD text helpers use the localization port rather than direct runtime string lookups.
\item
  DOGFOOD still keeps catalog loading and locale activation at the app composition boundary.
\end{itemize}

\section{Tests To Write First}\label{record-195-tests-to-write-first}

\begin{itemize}
\tightlist
\item
  Port resolution returns a localized object for a translated message.
\item
  Port resolution reports fallback status when selected-locale catalogs omit a value and fallback catalogs provide one.
\item
  Port resolution reports missing status and marker value when missing-marker policy is enabled.
\item
  Resource/data entries return localized object values without converting them to strings.
\item
  Localized objects and their issue/fact arrays are immutable.
\item
  DOGFOOD text lookup consumes \texttt{L\allowbreak{}o\allowbreak{}c\allowbreak{}a\allowbreak{}l\allowbreak{}i\allowbreak{}z\allowbreak{}a\allowbreak{}t\allowbreak{}i\allowbreak{}o\allowbreak{}n\allowbreak{}P\allowbreak{}o\allowbreak{}r\allowbreak{}t} and does not call \texttt{I\allowbreak{}1\allowbreak{}8\allowbreak{}n\allowbreak{}R\allowbreak{}u\allowbreak{}n\allowbreak{}t\allowbreak{}i\allowbreak{}m\allowbreak{}e\allowbreak{}.\allowbreak{}t\allowbreak{}(\allowbreak{})} directly.
\end{itemize}

\section{Retrospective}\label{record-195-retrospective}

Not started.

